\documentstyle[%


righttag%


,11pt%


,amssymb%


,russian%


,izvru%


]{amsart}





\newcommand{\tl}[3]{\medskip\noindent{\bf#1}\ #2 \dotfill #3 \par} 


\pagestyle{empty}


% \pagestyle{plain}


\begin{document}


\par


\vskip 2cm


\par


\bigskip


\bigskip


\bigskip


\par


\centerline{\Large СОДЕРЖАНИЕ}


\bigskip


\bigskip


%%%%%%%%%% Use \newline %%%%%%%%%%%





\tl{Винник А.В., Лейко С.Г.} 


{Изопериметрические экстремали поворота на двумерных \newline


связных группах Ли с инвариантными римановыми метриками}{3}


\tl{Горлов С.К., Новиков И.Я., Родин В.А.} 


{Коррекция полиномов Хаара, применяемых \newline


для сжатия графической информации}{6}


\tl{Игошин В.А.} 


{Пульверизационное моделирование и точечно-траекторные 


морфизмы\newline квазигеодезических потоков}{11}


\tl{Левенштам В.Б.} 


{О методе усреднения для 


квазилинейных параболических уравнений\newline с 


быстро осциллирующими коэффициентами}{22}


\tl{Лернер Э.Ю.} 


{Об адельной формуле для гауссовых интегралов}{31}


\tl{Марюкова Н.Е.} 


{Об одной геометрической интерпретации 


уравнения $\sin$-Гордона}{37}


\tl{Можей Н.П.}


{Однородные подмногообразия в четырехмерной 


аффинной и проективной\newline геометрии}{41}


\tl{Семигродских А.П.}


{О клонах полиномов над бесконечными полями}{53}


\tl{Сорокина М.М.} 


{О композиционных и локальных критических формациях}{59}


\tl{Шамоян Р.Ф.} 


{Непрерывные функционалы и мультипликаторы степенных рядов одного \newline


класса аналитических в поликруге функций}{67}





\bigskip





\centerline{КРАТКИЕ СООБЩЕНИЯ}


\bigskip





\tl{Аминова А.В., Зуев С.В., Калинин Д.А.}


{Алгебраические условия согласования\newline двух метрических форм с общей


почти комплексной (кватернионной) структурой на\newline многообразии}{70}


\tl{Гайсин Т.И.}


{О комплексе базовых форм слоения}{74}


\tl{Крючков А.Ф., Сиренек В.А.} 


{Вероятностное представление решения 


гиперболического\newline уравнения массопереноса 


с конвективным членом}{77}


\tl{Чилин В.И., Ганиев И.Г.} 


{Индивидуальная эргодическая теорема для сжатий  \newline


в решетке $L_p(\wh\nabla,\wh\mu)$ Банаха--Канторовича}{81}








\vfill


\noindent\raisebox{2pt}{\copyright} Казанский госудаpственный унивеpситет


\end{document}


